\documentclass{report}
\usepackage{luatexja} % LuaTeXで日本語を使うためのパッケージ
\usepackage{luatexja-fontspec} % LuaTeX用の日本語フォント設定

% --- 数学関連 ---
\usepackage{amsmath, amssymb, amsfonts, mathtools, bm, amsthm} % 基本的な数学パッケージ
\usepackage{type1cm, upgreek} % 数式フォントとギリシャ文字
\usepackage{physics, mhchem} % 物理や化学の記号や式の表記を簡単にする

% --- 表関連 ---
\usepackage{multirow, longtable, tabularx, array, colortbl, dcolumn, diagbox} % 表のレイアウトを柔軟にする
\usepackage{tablefootnote, truthtable} % 表中に注釈を追加、真理値表
\usepackage{tabularray} % 高度な表組みレイアウト

% --- グラフィック関連 ---
\usepackage{tikz, graphicx} % 図の描画と画像の挿入
\usepackage{background} % ウォーターマークの設定
\usepackage{caption, subcaption} % 図や表のキャプション設定
\usepackage{float, here} % 図や表の位置指定

% --- レイアウトとページ設定 ---
\usepackage{fancyhdr} % ページヘッダー、フッター、余白の設定
\usepackage[top = 20truemm, bottom = 20truemm, left = 20truemm, right = 20truemm]{geometry}
\usepackage{fancybox, ascmac} % ボックスのデザイン

% --- 色とスタイル ---
\usepackage{xcolor, color, colortbl, tcolorbox} % 色とカラーボックス
\usepackage{listings, jvlisting} % コードの色付けとフォーマット

% --- 参考文献関連 ---
\usepackage{biblatex, usebib} % 参考文献の管理と挿入
\usepackage{url, hyperref} % URLとリンクの設定

% --- その他の便利なパッケージ ---
\usepackage{footmisc} % 脚注のカスタマイズ
\usepackage{multicol} % 複数段組
\usepackage{comment} % コメントアウトの拡張
\usepackage{siunitx} % 単位の表記
\usepackage{docmute}
% \usepackage{appendix}
% --- tcolorboxとtikzの設定 ---
\tcbuselibrary{theorems, breakable} % 定理のボックスと改ページ設定
\usetikzlibrary{decorations.markings, arrows.meta, calc} % tikzの装飾や矢印の設定

% --- 定理スタイルと数式設定 ---
\theoremstyle{definition} % 定義スタイル
\numberwithin{equation}{section} % 式番号をサブセクション単位でリセット

% --- hyperrefの設定 ---
\hypersetup{
  setpagesize = false,
  bookmarks = true,
  bookmarksdepth = tocdepth,
  bookmarksnumbered = true,
  colorlinks = false,
  pdftitle = {}, % PDFタイトル
  pdfsubject = {}, % PDFサブジェクト
  pdfauthor = {}, % PDF作者
  pdfkeywords = {} % PDFキーワード
}

% --- siunitxの設定 ---
\sisetup{
  table-format = 1.5, % 小数点以下の桁数
  table-number-alignment = center, % 数値の中央揃え
}

% --- 透かし画像の設定 --- 
\backgroundsetup{
  scale=0.5,                       % 画像のスケール
  % color=black,                   % 画像の色（透かし用に半透明が推奨）
  opacity=0.2,                   % 透かしの透明度（0が完全透明、1が完全不透明）
  angle=0,                       % 画像の角度
  position = current page.south east,  % ページの右下
  hshift=-6cm, % 右方向へのシフト（負の値で内側に移動）
  vshift=5cm,  % 上方向へのシフト（正の値で内側に移動）
  contents={\includegraphics{./fig/appilogo-circular-full.png}} % 画像のパス
}

% --- その他の設定 ---
\allowdisplaybreaks % 数式の途中改ページ許可
\newcolumntype{t}{!{\vrule width 0.1pt}} % 新しいカラムタイプ
\newcolumntype{b}{!{\vrule width 1.5pt}} % 太いカラム
\UseTblrLibrary{amsmath, booktabs, counter, diagbox, functional, hook, html, nameref, siunitx, varwidth, zref} % tabularrayのライブラリ
\setlength{\columnseprule}{0.4pt} % カラム区切り線の太さ
\captionsetup[figure]{font = bf} % 図のキャプションの太字設定
\captionsetup[table]{font = bf} % 表のキャプションの太字設定
\captionsetup[lstlisting]{font = bf} % コードのキャプションの太字設定
\captionsetup[subfigure]{font = bf, labelformat = simple} % サブ図のキャプション設定
\setcounter{secnumdepth}{4} % セクションの深さ設定
\newcolumntype{d}{D{.}{.}{5}} % 数値のカラム
\newcolumntype{M}[1]{>{\centering\arraybackslash}m{#1}} % センター揃えのカラム
\DeclareMathOperator{\diag}{diag}
\everymath{\displaystyle} % 数式のスタイル
\newcommand{\inner}[2]{\left\langle #1, #2 \right\rangle}
\renewcommand{\figurename}{図}
\renewcommand{\i}{\mathrm{i}} % 複素数単位i
\renewcommand{\laplacian}{\grad^2} % ラプラシアンの記号
\renewcommand{\thesubfigure}{(\alph{subfigure})} % サブ図の番号形式
\newcommand{\m}[3]{\multicolumn{#1}{#2}{#3}} % マルチカラムのショートカット
\renewcommand{\r}[1]{\mathrm{#1}} % mathrmのショートカット
\newcommand{\e}{\mathrm{e}} % 自然対数の底e
\newcommand{\Ef}{E_{\mathrm{F}}} % フェルミエネルギー
\renewcommand{\c}{\si{\degreeCelsius}} % 摂氏記号
\renewcommand{\d}{\r{d}} % d記号
\renewcommand{\t}[1]{\texttt{#1}} % タイプライタフォント
\newcommand{\kb}{k_{\mathrm{B}}} % ボルツマン定数
% \renewcommand{\phi}{\varphi} % ϕをφに変更
\renewcommand{\epsilon}{\varepsilon}
\newcommand{\fullref}[1]{\textbf{\ref{#1} \nameref{#1}}}
\newcommand{\reff}[1]{\textbf{図\ref{#1}}} % 図参照のショートカット
\newcommand{\reft}[1]{\textbf{表\ref{#1}}} % 表参照のショートカット
\newcommand{\refe}[1]{\textbf{式\eqref{#1}}} % 式参照のショートカット
\newcommand{\refp}[1]{\textbf{コード\ref{#1}}} % コード参照のショートカット
\renewcommand{\lstlistingname}{コード} % コードリストの名前
\renewcommand{\theequation}{\thesection.\arabic{equation}} % 式番号の形式
\renewcommand{\footrulewidth}{0.4pt} % フッターの線
\newcommand{\mar}[1]{\textcircled{\scriptsize #1}} % 丸囲み文字
\newcommand{\combination}[2]{{}_{#1} \mathrm{C}_{#2}} % 組み合わせ
\newcommand{\thline}{\noalign{\hrule height 0.1pt}} % 細い横線
\newcommand{\bhline}{\noalign{\hrule height 1.5pt}} % 太い横線

% --- カスタム色定義 ---
\definecolor{burgundy}{rgb}{0.5, 0.0, 0.13} % バーガンディ色
\definecolor{charcoal}{rgb}{0.21, 0.27, 0.31} % チャコール色
\definecolor{forest}{rgb}{0.0, 0.35, 0} % 森の緑色

% --- カスタム定理環境の定義 ---
\newtcbtheorem[number within = chapter]{myexc}{練習問題}{
  fonttitle = \gtfamily\sffamily\bfseries\upshape,
  colframe = forest,
  colback = forest!2!white,
  rightrule = 1pt,
  leftrule = 1pt,
  bottomrule = 2pt,
  colbacktitle = forest,
  theorem style = standard,
  breakable,
  arc = 0pt,
}{exc-ref}
\newtcbtheorem[number within = chapter]{myprop}{命題}{
  fonttitle = \gtfamily\sffamily\bfseries\upshape,
  colframe = blue!50!black,
  colback = blue!50!black!2!white,
  rightrule = 1pt,
  leftrule = 1pt,
  bottomrule = 2pt,
  colbacktitle = blue!50!black,
  theorem style = standard,
  breakable,
  arc = 0pt
}{proposition-ref}
\newtcbtheorem[number within = chapter]{myrem}{注意}{
  fonttitle = \gtfamily\sffamily\bfseries\upshape,
  colframe = yellow!20!black,
  colback = yellow!50,
  rightrule = 1pt,
  leftrule = 1pt,
  bottomrule = 2pt,
  colbacktitle = yellow!20!black,
  theorem style = standard,
  breakable,
  arc = 0pt
}{remark-ref}
\newtcbtheorem[number within = chapter]{myex}{例題}{
  fonttitle = \gtfamily\sffamily\bfseries\upshape,
  colframe = black,
  colback = white,
  rightrule = 1pt,
  leftrule = 1pt,
  bottomrule = 2pt,
  colbacktitle = black,
  theorem style = standard,
  breakable,
  arc = 0pt
}{example-ref}
\newtcbtheorem[number within = chapter]{exc}{Requirement}{myexc}{exc-ref}
\newcommand{\rqref}[1]{{\bfseries\sffamily 練習問題 \ref{exc-ref:#1}}}
\newtcbtheorem[number within = chapter]{definition}{Definition}{mydef}{definition-ref}
\newcommand{\dfref}[1]{{\bfseries\sffamily 定義 \ref{definition-ref:#1}}}
\newtcbtheorem[number within = chapter]{prop}{命題}{myprop}{proposition-ref}
\newcommand{\prref}[1]{{\bfseries\sffamily 命題 \ref{proposition-ref:#1}}}
\newtcbtheorem[number within = chapter]{rem}{注意}{myrem}{remark-ref}
\newcommand{\rmref}[1]{{\bfseries\sffamily 注意 \ref{remark-ref:#1}}}
\newtcbtheorem[number within = chapter]{ex}{例題}{myex}{example-ref}
\newcommand{\exref}[1]{{\bfseries\sffamily 例題 \ref{example-ref:#1}}}
% --- 再定義コマンド ---
% \mathtoolsset{showonlyrefs=true} % 必要な式番号のみ表示
\pagestyle{fancy} % ヘッダー・フッターのスタイル設定
\chead{応用量子物性講義ノート} % 中央ヘッダー
% \rhead{}
\fancyhead[R]{\rightmark}
\renewcommand{\sectionmark}[1]{\markright{\thesection\ #1}}
\cfoot{\thepage} % 中央フッターにページ番号
\lhead{}
\rfoot{Yuto Masuda and Haruki Aoki} % 右フッターに名前
\setcounter{tocdepth}{4} % 目次の深さ
\makeatletter
\@addtoreset{equation}{section} % サブセクションごとに式番号をリセット
\makeatother

% --- メタ情報 ---
\title{応用量子物性講義ノート}
\date{更新日\today}
\author{Yuto Masuda and Haruki Aoki}

\begin{document}
  次に局所ゲージ変換を議論する．
  波動関数$\psi(\bm{r}, t)$を，
  \begin{align}
    \psi(\bm{r}, t) \to \psi'(\bm{r}, t) \coloneqq \e^{\i\alpha(\bm{r}), t}\psi(\bm{r}, t)
  \end{align}
  のように変換することを局所ゲージ変換という．
  局所ゲージ変換では，$\alpha$が$\bm{r}$に依存することに注意する．
  簡単な考察により，局所ゲージ変換にSchr\"odinger方程式は影響を受けることが分かる．
  \par
  例えば波動関数の勾配を計算すると，
  \begin{align}
    \grad\psi'(\bm{r}, t) &= \psi(\bm{r})\grad\e^{\i\alpha(\bm{r}, t)} + \e^{\i\alpha(\bm{r}, t)}\grad\qty(\psi(\bm{r}, t)) \\ 
    &= \i\qty(\grad\alpha(\bm{r}, t))  \e^{i\alpha(\bm{r}, t)}\psi(\bm{r}) + \e^{\i\alpha(\bm{r}, t)}\grad\psi(\bm{r}, t) \\ 
    &= \e^{\i\alpha(\bm{r}, t)}(\grad + \i\grad\alpha(\bm{r}, t))\psi(\bm{r}, t) 
  \end{align}
  となり，$\grad\psi$と$\grad\psi'$が一致しない．
  時間発展するSchr\"odinger方程式と物理量$\hat{A}$の期待値についても同様に計算を行えば，
  \begin{align}
    \begin{dcases}
      \i\hbar\pdv{t}\psi' \neq - \frac{\hbar^2}{2m}\laplacian\psi' \\
      \mel**{\psi}{\hat{A}}{\psi} \neq \mel**{\psi'}{\hat{A}}{\psi'}
    \end{dcases}
  \end{align}
  である．したがって，局所ゲージ変換に対して物理は不変ではない．そこで，\textbf{局所ゲージ不変性を基本原理とする物理を再構築する．}% モチベ？？相対論？
  \par
  今まで用いていた物理の模型の欠点は，積の微分をおこなったときに余分な項が出てきてしまうことであった．
  そこで，局所ゲージ不変性を満たすように物理の模型を変更するために，微分の定義を変更することを考える．
  以下では，今まで考えていた模型において空間に対する微分であった勾配を，共変微分$\hat{\bm{D}}$で，時間微分を$\hat{D}_t$で書き換えることよって，局所ゲージ不変な物理模型を作る．
  $\hat{\bm{D}}$と$\hat{D}_t$を，
  \begin{align}
    \hat{\bm{D}} &\coloneqq \grad + \i\frac{e}{\hbar}\hat{\bm{A}} \label{hatd-def}\\ 
    \hat{D}_t &\coloneqq \pdv{t} - \i\frac{e}{\hbar}\phi \label{hatdt-def}
  \end{align}
  と定義する．
  また，$\hat{\bm{A}}$と$\phi$の局所ゲージ変換則を，
  \begin{align}
    \hat{\bm{A}} &\to \hat{\bm{A}} - \frac{\hbar}{e}\grad\alpha(\bm{r}, t) \eqqcolon \hat{\bm{A}}' \\ 
    \phi &\to \phi + \frac{\hbar}{e}\pdv{t}\alpha(\bm{r}, t) \eqqcolon \phi'  
  \end{align}
  と定義する．$\hat{\bm{D}}$と$\hat{D}_t$を局所ゲージ変換すると，
  \begin{align}
    \hat{\bm{D}} &\to \hat{\bm{D}}' = \grad + \i\frac{e}{\hbar}\hat{\bm{A}'} = \grad + \i\frac{e}{\hbar}\bm{\hat{A}} - \i\grad\alpha(\bm{r}, t) \\ 
    \hat{D}_t &\to \hat{D}_t' = \pdv{t} - \i\frac{e}{\hbar}\phi' = \pdv{t} - \i\frac{e}{\hbar}\phi - \i\pdv{t}\alpha(\bm{r}, t)
  \end{align}
  となる．
  このように微分を定義すると，
  \begin{align}
    \hat{\bm{D}}'\psi'(\bm{r}, t) &= \e^{\i\alpha(\bm{r}, t)}\hat{\bm{D}}\psi(\bm{r}, t) \label{d-prime}\\ 
    \hat{\bm{D}}'^2\psi'(\bm{r}, t) &= \e^{\i\alpha(\bm{r}, t)}\hat{\bm{D}}^2\psi(\bm{r}, t) \label{d-prime-power}\\ 
    \hat{D}_t'\psi'(\bm{r}, t) &= \e^{\i\alpha(\bm{r}, t)}\hat{D}_t\psi(\bm{r}, t) \label{dt-prime}
  \end{align}
  となることが確かめられる．
  \par
  \refe{d-prime-power}や\refe{dt-prime}を使いながら，元のSchr\"odinger方程式において，勾配を，$\hat{\bm{D}}$で，時間微分を$\hat{D}_t$で書き換えると，
  \begin{align}
    \i\hbar D'_t\psi' &= -\frac{\hbar^2}{2m}\bm{D}'^2\psi' \\
    \Leftrightarrow \e^{\i\alpha(\bm{r}, t)}\i\hbar D_t\psi &= -\e^{\i\alpha(\bm{r}, t)}\frac{\hbar^2}{2m}\bm{D}^2\psi\label{local-gauge-psi}
  \end{align}
  となり，局所ゲージ変換に対して不変なSchrödinger方程式となった．
  また，\refe{local-gauge-psi}において\refe{hatd-def}と\refe{hatdt-def}を\refe{local-gauge-psi}に代入すると，
  \begin{itembox}[l]{局所ゲージ変換に対して不変なSchrödinger方程式}
    \begin{align}
      \i\hbar\pdv{t}\psi = \qty[\frac{1}{2m}\qty(\hat{\bm{p}} + e\hat{\bm{A}})^2 - e\phi]\psi
    \end{align}
  \end{itembox}
  を得る．
  \par
  以上の流れをまとめると，局所ゲージ不変性を要請することにより，ゲージ場$\bm{A}$が導入された．逆に，電磁場の起源は局所ゲージ不変性であるといえる．

  \begin{myexc}{電磁場中のSchrödinger方程式}{}
  \begin{enumerate}
    \item 共変微分
      \begin{align}
        \hat{\bm{D}} = \nabla + \i\frac{e}{\hbar}\bm{A}
      \end{align}
      において，
      \begin{align}
        \hat{\bm{D}} \times \hat{\bm{D}} = \i\frac{e}{\hbar}\bm{B}
      \end{align}
      が成り立つことを示せ．
    \item 電磁場中のSchrödinger方程式
    \begin{align}
      \i\hbar\pdv{t}\psi = \qty[\frac{1}{2m}\qty(\hat{\bm{p}} + e\hat{\bm{A}})^2 - e\phi]\psi
    \end{align}
    においては
    \begin{align}
      \frac{\partial}{\partial t} \rho + \nabla\cdot \bm{j} = 0 \label{prob-current-exercise}
    \end{align}
    が成り立たない．ここで，
    \begin{align}
      \bm{j} = \frac{\i\hbar}{2m}\qty[(\nabla \psi^{*})\psi - \psi^{*} (\nabla\psi)],\ \rho = \psi^{*}\psi
    \end{align}
    である．\refe{prob-current-exercise}が成り立つように確率の流れを再定義せよ．
  \end{enumerate}
  \tcblower
  \begin{enumerate}
    \item \begin{align}
      (\hat{\bm{D}} \times \hat{\bm{D}})\psi &= \qty(\nabla + \i\frac{e}{\hbar}\bm{A}) \times \qty(\nabla + \i\frac{e}{\hbar}\bm{A})\psi\\
      &= \qty(\nabla\times\nabla + \i\frac{e}{\hbar}\bm{A}\times\nabla + \i\frac{e}{\hbar}\nabla\times\bm{A} - \frac{e^2}{\hbar^2})\psi\\
      &= \i\frac{e}{\hbar}\bm{A}\times(\nabla\psi) + \i\frac{e}{\hbar}(\nabla\psi)\times\bm{A} + \i\frac{e}{\hbar}(\nabla\times\bm{A})\psi\\
      &= \bm{B}\psi
    \end{align}
    \item まず，確率密度の時間微分は
      \begin{align}
        \frac{\partial}{\partial t}\rho = \qty(\frac{\partial}{\partial t}\psi^{*}) + \psi^{*}\qty(\frac{\partial}{\partial t}\psi)
      \end{align}
      である．次に，電磁場中のSchrödinger方程式より，
      \begin{align}
        \i\hbar \frac{\partial}{\partial t} \psi &= -\frac{\hbar^2}{2m}\qty(\nabla + \i\frac{e}{\hbar}\bm{A})^2\psi - e\phi\psi\\
        \i\hbar \frac{\partial}{\partial t} \psi^{*} &= \frac{\hbar^2}{2m}\qty(\nabla - \i\frac{e}{\hbar}\bm{A})^2\psi^{*} - e\phi\psi^{*}
      \end{align}
      を得る．ここで，
      \begin{align}
        \qty(\nabla - \i\frac{e}{\hbar}\bm{A})^2\psi^{*} &= \nabla^2\psi^{*} - \i\frac{e}{\hbar}\bm{A}\cdot(\nabla\psi^{*}) - \i\frac{e}{\hbar}(\nabla\psi^{*})\cdot\bm{A} - \i\frac{e}{\hbar}(\nabla\cdot\bm{A})\psi^{*} - \frac{e^2}{\hbar^2}\psi^{*}\\
        &= \nabla^2\psi^{*} - 2\i\frac{e}{\hbar}\bm{A}\cdot(\nabla\psi^{*}) - \i\frac{e}{\hbar}(\nabla\cdot\bm{A})\psi^{*} - \frac{e^2}{\hbar^2}\psi^{*}\\
        \qty(\nabla + \i\frac{e}{\hbar}\bm{A})^2\psi &=\nabla^2\psi + 2\i\frac{e}{\hbar}\bm{A}\cdot(\nabla\psi) + \i\frac{e}{\hbar}(\nabla\cdot\bm{A})\psi - \frac{e^2}{\hbar^2}\psi
      \end{align}
      である．よって，
      \begin{align}
        2m\i\frac{\partial}{\partial t}\rho &= -\nabla\cdot\qty[\psi^{*}\nabla\psi - (\nabla\psi^{*})\psi - 2\i \frac{e}{\hbar}\bm{A}\psi^{*}\psi]\\
        &= -\nabla\cdot\qty[\psi^{*}(\nabla - \i\frac{e}{\hbar} \bm{A})\psi - ((\nabla - \i \frac{e}{\hbar}\bm{A})\psi)^{*}\psi]
      \end{align}
      となるため，
      \begin{align}
        \bm{j} = \frac{\i}{2m}\qty[\psi^{*}(\bm{D}\psi) - (\bm{D}^{*}\psi^{*})\psi]
      \end{align}
      と定義すればよいことがわかる．
  \end{enumerate}
  \end{myexc}
\end{document}